% Source: https://artofproblemsolving.com/community/c3h2083733_a_magical_coin_problem
% Stable identity: aops-topic-2083733-q01
Let \(p_n\) be the probability that flip \(n\) is heads, so \(p_1=\dfrac{1}{2}\). After heads, the next flip is heads with probability \(a=\dfrac{1}{2}\left(1+\dfrac{h}{100}\right)\). After tails, it is tails with probability \(\dfrac{1}{2}\left(1+\dfrac{t}{100}\right)\), so heads with probability \(c=\dfrac{1}{2}\left(1-\dfrac{t}{100}\right)\). By the law of total probability,
\[
p_{n+1}=a\,p_n+c\,(1-p_n).
\]
\begin{enumerate}
\item[(a)] Here \(a=\dfrac{3}{4}\), \(c=\dfrac{1}{4}\), so \(p_{n+1}=\dfrac{1}{4}+\dfrac{1}{2}p_n\). Starting from \(p_1=\dfrac{1}{2}\) gives \(p_2=\dfrac{1}{2}\), and so on: every \(p_n=\dfrac{1}{2}\) (as expected by the heads-tails symmetry). The answer is \(\boxed{\dfrac{1}{2}}\).
\item[(b)] Here \(a=\dfrac{1}{2}\cdot\dfrac{7}{4}=\dfrac{7}{8}\) and \(c=\dfrac{1}{2}\cdot\dfrac{3}{4}=\dfrac{3}{8}\), so \(p_{n+1}=\dfrac{3}{8}+\dfrac{1}{2}p_n\). Then
\[
p_2=\dfrac{5}{8},\quad p_3=\dfrac{11}{16},\quad p_4=\dfrac{23}{32},\quad p_5=\dfrac{47}{64}.
\]
The answer is \(\boxed{\dfrac{47}{64}}\).
\end{enumerate}
